\chapter{Model Configurations and Implementation Details}
\label{app:implementation}

\section{Executed Component Boundaries}
The PHM scalar-image regressor and the REASSEMBLE action classifier use different
components of RT-DETR. Table~\ref{tab:component-boundaries} identifies the model
weights, intermediate representation and output used by each branch.

\begin{table}[htbp]
\caption{Inputs and Outputs of the Executed Model Components}\label{tab:component-boundaries}
\centering\small

\begin{tabular}{P{0.19\linewidth}P{0.34\linewidth}P{0.37\linewidth}}\hline
Component & Input and representation & Executed boundary\\\hline
PHM scalar image branch & Three spatially pooled 256-dimensional maps, plus recorded
extraction latency: 769 actual head covariates & Ultralytics RT-DETR-L prediction path
with hooked maps; not the CNN-only REASSEMBLE path\\
REASSEMBLE visual branch & Sixteen RGB frames pooled to an 896-dimensional action
descriptor & Frozen ResNet18-vd CNN backbone inside RT-DETR-R18vd, followed by a trained
binary head\\
REASSEMBLE numerical statistics & 22 channels, ten statistics each: 220 features &
Training-only standardization and logistic regression\\
REASSEMBLE temporal branch & 512 progress positions by 22 channels & Scratch-initialized
compact PatchTST and fitted projection/head\\
External numerical adapters & MSM: 90; HANDOVER: 36; ViFailback: 280 statistics &
Dataset-specific input construction; not an unchanged raw-sequence temporal
encoder\\\hline
\end{tabular}
\displaynote{The reported dimensions apply to the implementation paths inspected.}
\end{table}
\ev{Executed component map}{\EvidenceRef{E110}{}}

The PHM scalar-image head receives extraction latency alongside pooled image features.
The dependence of latency on the computing environment limits attribution of the
regressor's performance to image content. The target measures candidate area in
photographs. Earlier target versions selected the maximum across views, while the
retained version uses a fixed view for each tooth. The top-three run summary aggregates
the tooth values. Independent physical defect masks were unavailable to validate
localization.

\section{Temporal Representations}
The PHM implementation defines a compact channel-independent PatchTST-style regressor
with default patch length sixteen, stride eight, hidden width 32, four attention heads,
two encoder layers, feed-forward width 64, dropout 0.1 and a 64-unit head. The
configured input width remains run-specific. A scalar output and the optional
sorted-profile head answer different target questions. The latter uses a base and
positive softplus increments to enforce ascending order; this guarantees an
order-statistic profile but discards tooth identity. The monotonicity does not
demonstrate that the estimated profile matches actual damage.\ev{PHM temporal
configuration}{\EvidenceRef{E031}{}}

The REASSEMBLE temporal branch follows a separate configuration: patch length 32, stride
sixteen, width 32, two layers, four heads, feed-forward width 64 and dropout 0.1. It
begins without pretrained temporal weights. Mean patch pooling and channel concatenation
feed a 128-dimensional projection and binary head, yielding 109,505 parameters in the
inspected implementation. Observed-value masks are retained and passed through the
relevant interface, but that does not establish attention suppression of every missing
token. Missing values are filled under training-only normalization. The source-specific
mask semantics therefore remain part of the implementation limitation rather than being
inferred from the model family \citep{Nie2023PatchTST}.\ev{REASSEMBLE temporal
recipe}{\EvidenceRef{E053}{}}

\section{Heads, Fitting and Calibration}
The original visual head is a new 896--128--1 network with ReLU and dropout 0.1,
containing 114,945 trainable parameters above a frozen backbone. Original visual epoch
candidates are five and ten; temporal candidates are three and six. Both use AdamW with
learning rate 0.001, weight decay 0.001 and batch size 128. These candidates constitute
the recorded selection procedure, not a broad optimizer search. Logistic regression
separately compares $C\in\{0.1,1,10\}$ under inner-fold selection.\ev{Original branch
configuration}{\EvidenceRef{E053}{}}

Branch calibration used an unweighted sigmoid logistic fit to the selected inner
out-of-fold scores. The outer assessment labels were withheld from calibration. The
original procedure selected a threshold between 0.01 and 0.99 by maximizing inner
out-of-fold balanced accuracy, resolving ties toward 0.5. Later studies used a fixed
threshold of 0.5. Comparisons between those procedures therefore combine the effects of
model and threshold changes.

\section{Fusion Mechanisms and Attribution}
Uniform fusion averages the two calibrated branch probabilities. F2 fits an unweighted
L2 logistic model to those probabilities, with $C=1$ and a new intercept and two
coefficients. It is neither an F-beta metric nor a convex weighted average. Its
coefficients depend on score scaling and are not causal percentages of modality
importance.

Feature fusion F3 combines a 128-dimensional visual projection with the 220 statistical
numerical features through separate projections, normalization and a fitted classifier.
Its recorded candidates use hidden widths 32 or 64 and ten or twenty epochs. Context
gates instead allocate weights over branch logits. Their candidates have hidden width
eight or sixteen and forty epochs. Action-only, quality-only and combined contexts
remain different models, as do a static gate with context held constant and a gate
exposed to corrupted OOF records.\ev{Fusion and gate
configurations}{\EvidenceRef{E054}{}}

With one branch available, the explicitly masked convex-logit gate gives that branch a
weight of one. The resulting prediction follows from the specified masking rule. When
both branches are absent, some implementations reject the input and the external
evaluation records abstention and coverage.

The later anchored adjustment is another distinct mechanism. It modifies frozen F2
probability coefficients by bounded context-dependent factors, with a zero-initialized
adaptive output. It is not the same convex logit gate and need not have the same
calibration or threshold policy. Appropriate static and exposure-matched controls are
necessary to identify what, if anything, adaptation adds.\ev{Anchored
implementation}{\EvidenceRef{E063}{}}

\section{External Adaptations}
MSM and HANDOVER use the shared frozen visual backbone with a newly trained head and a
dataset-specific numerical classifier. ViFailback adds pooling across three camera
views. Its matched study compares frozen generic and transferred CNN weights with a
fresh target head in each case. Scalers, sigmoid calibrators and out-of-fold fusion
coefficients were fitted to target data where specified. The separate condition for
updating all target CNN weights was incomplete at the November assessment.\ev{External
model and completion register}{\EvidenceRef{E108}{}}
